% !TEX program = xelatex
\documentclass[11pt]{article}
\usepackage[margin=1in]{geometry}
\usepackage{parskip}
\usepackage{hyperref}
\pagestyle{empty}

\begin{document}

\begin{flushright}
Junjie Chen\\
School of Computer Science and Engineering\\
Nanjing University of Science and Technology\\
Nanjing, China\\
\href{mailto:junjie_chen@njust.edu.cn}{junjie\_chen@njust.edu.cn}
\end{flushright}

\vspace{1em}
\noindent 17 October 2026

\vspace{1em}
\noindent Dear Editor-in-Chief,

\vspace{0.5em}
Please consider the enclosed manuscript, \emph{Identifiability and Confounding in
Test-Adequacy Metric Validity Studies: A Cross-Language Empirical Analysis}, for
publication in the \emph{Journal of Systems and Software} as a research paper.

The manuscript examines the evidential basis of a standard argument in empirical
software engineering: that a test-adequacy metric is incrementally valuable
because adding it to a regression improves the cross-validated $R^2$. Rather than
proposing a new metric, we ask whether that argument can be statistically
resolved at the sample sizes at which it is routinely made. Using Defects4J
(Java) and BugsInPy (Python), we show that the increments being claimed are of
the same order as their own uncertainty; that their sign depends on
cross-validation and resampling choices that are rarely reported; and that
project-level confounding alone can manufacture a material increment in up to
82\% of runs, a rate that project fixed effects hold below 1\%.

Our contribution is methodological rather than a benchmark result, and we
believe it speaks to a broad part of the journal's readership: the same
incremental-prediction design is used to validate metrics, defect predictors,
and effort models across software engineering. We close with a five-check
protocol that requires no test executions and can be applied to any published
increment-in-$R^2$ claim, including those in already-published work.

All scripts, collected measurements, and analysis outputs are publicly archived
(DOI: \href{https://doi.org/10.5281/zenodo.23432059}{10.5281/zenodo.23432059})
and the full pipeline is available at
\url{https://github.com/JunjieChen0/test-adequacy-resolution}. Every number in
the manuscript is produced by the shipped analysis code from the shipped
measurements; none is entered by hand.

The manuscript is original, has not been published previously, and is not under
consideration by another journal. As the sole author, I have approved the
submission, declare no competing interests, and take full responsibility for
the content. A declaration of generative AI and AI-assisted technologies is
included in the manuscript.

\vspace{1em}
\noindent Thank you for considering this submission.

\vspace{1.5em}
\noindent Sincerely,\\[0.3em]
Junjie Chen

\end{document}
